\section{The target and the meaning of ``simpler''}
\label{sec:overview}

Let $\Con(\ZFC)$ mean that there is no finite derivation of a contradiction
from the recursively presented first-order ZFC axioms, in the proof
calculus used by the repository's K215 calibration.  Equivalent standard
effective proof calculi give provably equivalent consistency statements.
The required conclusion has the following form for one fully fixed
$F\in\Z[X_1,\ldots,X_n]$:
\begin{equation}\label{eq:target}
 \ZFC\ \vdash\quad
 \bigl[\neg\exists\boldsymbol{x}\in\Z^n\ F(\boldsymbol{x})=0\bigr]
 \ \longleftrightarrow\ \Con(\ZFC).
\end{equation}
This is a proof of a biconditional.  It does not prove that either side
holds.  Under the usual consistency hypothesis, G\"odel's second
incompleteness theorem precludes ZFC from proving the left side once
the biconditional has been established.

There are several inequivalent notions of simplicity.  Fewer variables
can cost enormous degree; a low-degree system can require more witnesses;
small arithmetic circuits can contain very large fixed integers.  A
positive-variable equation also differs from an equation in unrestricted
integers.  We therefore retain several useful points instead of calling
one construction an absolute minimum.

\subsection{Concrete improvements over the pinned collection}

All comparisons are against the immutable snapshot in \cite{repo}, not
against a moving branch.  ``Operation'' means a binary addition,
subtraction or multiplication in the emitted directed acyclic graph.
Sharing is allowed; all powers of variables are expanded into such
gates.  Fixed integers are leaves.  A system is evaluated as shared left
and right sides to be compared; the comparisons themselves are not
arithmetic gates.  A sum of squared residuals includes the subtraction,
squaring and summation gates.  These conventions match the compared
repository receipts.

\begin{center}
\small
\begin{tabular}{@{}p{43mm}rrrr@{}}
\toprule
Form & Domain & Variables & Degree & Operations\\
\midrule
New Sun/Jones synthesis & $\Z$ & 10 & $d_{10}$ & 1,442\\
Jones wrapper, package 29 & $\Z$ & 17 & $d_{17}$ & 20,607\\
New Jones wrapper & $\Z$ & 16 & $d_{16}$ & 1,617\\
Quadratic system, package 30 & $\Z_{>0}$ & 31 & 2 & 65\\
New quadratic system & $\Z_{>0}$ & 31 & 2 & 60\\
Quartic, package 30 & $\Z_{>0}$ & 31 & 4 & 130\\
New quartic & $\Z_{>0}$ & 31 & 4 & 125\\
Product, package 30 & $\Z_{>0}$ & 12 & 1,081 & 111\\
New product evaluator & $\Z_{>0}$ & 12 & 1,081 & 108\\
System, package 30 & $\Z_{>0}$ & 12 & 1,072 & 107\\
New shifted system & $\Z_{>0}$ & 12 & 1,072 & 102\\
\bottomrule
\end{tabular}
\end{center}
Here
\[
 d_{10}=126496\,5^{31}+27064,\qquad
\]
\[
 d_{17}=424296\,5^{33}+53824,\qquad
 d_{16}=813056\,5^{32}+10016.
\]
The new ten-integer equation is specified in full in
Section~\ref{sec:signed10}.  Its exact degree is
\[
589042901992797851562527064,
\]
and its evaluator uses 851 multiplications
and 591 additions/subtractions.  It combines the fixed quartic input
with Sun's ten-integer reduction \cite{sun2021,sun2026} and the local
simplifications proved there.  The contribution is an explicit
$\Con(\ZFC)$ instance and smaller local formulas, not a new general
ten-unknowns theorem.

The collection already contains nine-variable equations over
$\N$ or $\Z_{>0}$.  Ten unrestricted integer variables should not be
compared with these as if the domains were identical.  The existing
67-operation single equation and the 59-operation higher-degree system
also remain better at their respective operation-count objectives.
Our 60-operation system improves the quadratic/quartic point.

\subsection{Fixed coefficients and the proof chain}

All the new forms are tied to one complete native coefficient tuple.
For its four ports the largest bit length falls from 1,958,376 to
1,070,004, while their combined raw storage falls from 977,220 to
533,099 bytes.  These are measurements of the \emph{native ports}; they
are not a claim about the largest coefficient after expanding an entire
high-degree polynomial.  In particular, the sixteen-variable Jones
wrapper has a much larger fixed base coefficient specified by a bounded
power-of-two recipe.

The logical composition is
\begin{align*}
 \text{a zero of the displayed equation/system}
 &\ \longleftrightarrow\ \text{a positive zero of the native parent}\\
 &\ \longleftrightarrow\ \text{a finite blank-input K215 halt}\\
 &\ \longleftrightarrow\ \text{a finite ZFC refutation}.
\end{align*}
Negating this equivalence yields \eqref{eq:target}.  The parent-to-machine
direction uses the native compiler and arithmetic-skeleton proofs in
the pinned repository.  The last equivalence is calibrated by
\texttt{blank\_halts\_iff\_zfc\_refutation} and
\texttt{nonhalting\_iff\_zfc\_consistency} in
\texttt{KEndpoint215.lean}.  The replacement tuple is justified by the
ordinary adapter proof in Section~\ref{sec:new-coefficients}.

No coefficient depends on deciding whether K215 halts.  The 430 machine
transitions, 943 tiles, masks, weighted quartic coefficients and all fixed
powers are bounded, explicitly reconstructible data.  Auxiliary names
inside the displayed formulas denote polynomial expressions, not
additional existential variables.  Fixed sums with large numerical
limits are ordinary polynomials; their efficient evaluators are supplied.

\subsection{Verification boundary}

The archive separates three kinds of evidence.  First are the inherited
semantic results, pinned source proofs and the existing Lean endpoint.
Second are the new ordinary mathematical proofs: the coefficient adapter,
square and divisibility mergers, the cofinal Jones restriction, and the
algebraic refactorings.  Third are executable finite checks of coefficient
bytes, symbolic identities, degree witnesses and circuit transcription.
The finite checks are useful for catching errors in a concrete equation;
they do not prove universal compiler correctness by testing examples.

The new work is not an end-to-end Lean or Rocq development.  In particular,
the ten-variable route imports a theorem from the explicitly identified
2026 preprint \cite{sun2026}; the sixteen-variable route independently
uses the Jones reduction and does not depend on that preprint.  The
ordinary proofs use finite arithmetic, induction, elementary algebra and
the cited Diophantine theorems, all available within ZFC.  No truth or
consistency assumption about ZFC is needed to prove the biconditional.
